% Avsnitt 17.1-17.8, ur lectures/L17/appendix/a_can_controller_receive.md.

\section{Att läsa tillbaka bussen}
\label{c17:sec:readback}

En mottagande nod samlar in riktiga (avstoppade) bitar genom \code{rx_shift_reg}, grindade vid
\code{bit_timer.sample}-punkten (mitt i perioden, när bussen hunnit stabilisera sig), inte vid
\code{bit_done}. Den sätter \code{rx_shift_reg.bit_count} till samma gruppstorlekar som sändarsidan
använder, och läser varje färdig grupp ur \code{rx_shift_reg.data}, som är \textbf{högerjusterad} i
sina lägsta \code{bit_count} bitar (\chapref{c15:ch}). Den sätter ihop ramen igen i ackumulatorer:
\begin{itemize}
  \item \textbf{ID}: de höga 8 bitarna från \code{ARB_HI} (\code{rxsr_data(7 downto 0)}), de låga 3
    från \code{ARB_LO} (\code{rxsr_data(3 downto 1)}). Gruppens sista bit, \code{rxsr_data(0)}, är
    RTR: dominant markerar de dataramar den här boken stöder, så en recessiv RTR höjer \code{error}
    och avbryter ramen (en fjärram, som inte stöds).
  \item \textbf{DLC}: från \code{CTRL} (\code{rxsr_data(3 downto 0)}). Gruppens två översta bitar,
    IDE och r0, kontrolleras \textbf{inte}: båda är alltid dominanta i den här bokens ramar. En
    mottagare som mötte den recessiva IDE:n i en riktig ram med utökad identifierare skulle feltolka
    den som en standardram ända tills CRC-kontrollen misslyckades; att kontrollera IDE och hantera
    det den annonserar hör hemma i en produktionskontroller (\secref{c19:sec:gaps} listar det bland
    luckorna).
  \item \textbf{Data}: varje mottagen byte placeras på sin \textbf{fasta} position efter byteindex,
    byte 0 i bitarna 63 downto 56, byte 1 i 55 downto 48, och så vidare. Det här är en positionsstyrd
    \code{case}, \textbf{inte} en löpande vänsterskiftning (som bara skulle landa byte 0 högst upp
    för en full 8-bytesram).
\end{itemize}

I slutet av \code{STATE_EOF}, lås \code{rx_id}, \code{rx_dlc} och \code{rx_data} från ackumulatorerna
och pulsa \code{rx_valid}.

% ------------------------------------------------------------------------------------------
\section{Mottagarvägen, tillstånd för tillstånd}
\label{c17:sec:states}

Samma tillståndssekvens som \chapref{c16:ch}:s sändväg, gången som mottagare. \code{role = '0'} hela
vägen. Läs det här vid sidan av \secref{c16:sec:fields}:s fälttabell: gruppstorlekarna är identiska,
bara det som händer med varje grupp skiljer sig. Tabellen är specifikationen;
\secref{c17:sec:complete} samlar de regler som inte ryms i den.

”Föregriper” är det \code{rxsr_bit_count}/\code{rxsr_enable} som det \textbf{kombinatoriska} blocket
driver medan \code{state} fortfarande läser \code{STATE_LOAD_*}-värdet, en cykel innan
skifttillståndet blir synligt; \secref{c17:sec:preempt} förklarar varför det är nödvändigt.

\begin{booktable}{@{}p{8.2em}p{4.6em}p{5.6em}L@{}}
  \thead{Tillstånd} & \thead{Föregriper} & \thead{Går vidare på} & \thead{Ackumulator / åtgärd} \\
  \midrule
  \code{STATE_IDLE} & inget & dominant \code{rx_bus_s2}, \code{idle_bits} full & lås
    \code{role = '0'}, nollställ \code{idle_bits} \\
  \code{STATE_START} & inget & omedelbart & nollställ \code{error}, återställ byteindexet,
    \code{bt_resync} \\
  \code{STATE_SOF} & 1 & \code{rxsr_valid} & samla in SOF som en enbitsgrupp; själva biten kastas \\
  \code{STATE_LOAD_ARB_HI} & 8 & \code{bt_sample} & (väntan: förra gruppens sista bit ligger
    fortfarande på ledningen) \\
  \code{STATE_ARB_HI} & 8 & \code{rxsr_valid} &
    \code{rx_id_acc(10 downto 3) <= rxsr_data(7 downto 0)} \\
  \code{STATE_LOAD_ARB_LO} & 4 & \code{bt_sample} & \\
  \code{STATE_ARB_LO} & 4 & \code{rxsr_valid} &
    \code{rx_id_acc(2 downto 0) <= rxsr_data(3 downto 1)}; \code{rxsr_data(0)} är RTR: om recessiv,
    \code{error <= '1'} och avbryt \\
  \code{STATE_LOAD_CTRL} & 6 & \code{bt_sample} & \\
  \code{STATE_CTRL} & 6 & \code{rxsr_valid} & \code{rx_dlc_acc <= rxsr_data(3 downto 0)};
    \code{DLC = 0} hoppar till \code{STATE_CRC_WAIT} \\
  \code{STATE_LOAD_DATA} & 8 & \code{bt_sample} & \\
  \code{STATE_DATA} & 8 & \code{rxsr_valid} & placera \code{rxsr_data} på den fasta positionen för
    byte \code{data_byte_idx} \\
  \code{STATE_CRC_WAIT} & inget & omedelbart & en stabiliseringscykel; \textbf{ingen} väntan på
    \code{bt_sample} \\
  \code{STATE_LOAD_CRC_HI} & 8 & \code{bt_sample} & \\
  \code{STATE_CRC_HI} & 8 & \code{rxsr_valid} & fortsätt mata \code{crc15}; själva värdet behålls
    inte \\
  \code{STATE_LOAD_CRC_LO} & 7 & \code{bt_sample} & \\
  \code{STATE_CRC_LO} & 7 & \code{rxsr_valid} & sista bitarna in i \code{crc15} \\
  \code{STATE_CRC_LO_WAIT} & inget & \code{bt_bit_done} & kontrollera \code{crc_valid}: om låg,
    \code{error <= '1'} och avbryt \\
  \code{STATE_CRC_DELIM} & inget & \code{bt_bit_done} & \\
  \code{STATE_ACK_SLOT} & inget & \code{bt_bit_done} & driv bussen dominant för att kvittera \\
  \code{STATE_ACK_DELIM} & inget & \code{bt_bit_done} & armera 7-bitars EOF-nedräkningen \\
  \code{STATE_EOF} & inget & \code{bt_bit_done} ×7 & vid den sista biten, lås \code{rx_id},
    \code{rx_dlc}, \code{rx_data}, pulsa \code{rx_valid} \\
\end{booktable}

Två mönster löper genom hela tabellen. Varje \code{STATE_LOAD_*}-rad finns enbart för att förbruka
den förra gruppens avslutande bitperiod: tillståndsmaskinen skördar ingenting där, även om det
\code{bt_sample} som avslutar tillståndet är precis det som \code{rx_shift_reg} samlar in nästa
grupps första bit på. Varje avslutande (ostoppad) rad går vidare på \code{bt_bit_done} i stället för
\code{bt_sample}, eftersom de fälten går förbi \code{rx_shift_reg} helt och räknas som hela
bitperioder.

\subsection*{Varför \code{STATE_SOF} aktiverar \code{rx_shift_reg} för en bit som ingen behåller}
Det insamlade värdet kastas; SOF är alltid dominant och \code{STATE_IDLE} visste redan det när den
startade den här ramen. Det som spelar roll är att biten passerar genom avstoppningens följdspårare.
Den spåraren (\code{consecutive}, \code{last_bit}) lever kvar över gruppgränser och nollställs bara
av reset (\chapref{c15:ch}), och sändarens kopia inuti \code{tx_shift_reg} räknade SOF, eftersom
\chapref{c16:ch} skickar den genom \code{tx_shift_reg} som vilket annat stoppat fält som helst. Hoppa
över den här, och mottagaren kör hela ramen en dominant bit efter sändaren.

Det är inte något litet fönster. Varje identifierare vars fyra översta bitar är noll, alltså
\textbf{varje ID under \code{0x080}}, ger SOF plus fyra dominanta identifierarbitar: sändarens följd
når fem och den sätter in en stoppbit, medan en mottagare som missade SOF bara står på fyra, inte
väntar sig någon, och tar den för en riktig bit. Allt efter den är förskjutet. Felet dyker upp långt
senare som ett CRC-fel eller ett stoppbitsbrott, och det slår mot precis de högprioriterade
identifierare som en riktig design reserverar för sina viktigaste meddelanden.

Regeln under det är värd att formulera för sig: \textbf{båda följdspårarna måste se exakt samma
bitar, SOF till och med slutet av CRC-fältet, och inget annat.} Svansen utesluts på båda sidor
eftersom inget av skiftregistren är aktiverat för den. Över \emph{ramgränser} är däremot ”samma
bitar” ingen garanti som designen kan luta sig mot alls, vilket är \secref{c17:sec:cleanslate}.

\code{STATE_SOF} behöver inget föregripande för det här. \code{STATE_START} pulsar \code{bt_resync},
så SOF-bitens sampelpunkt ligger ett helt \code{SAMPLE_TICK} framåt (\chapref{c12:ch});
tillståndsmaskinen är med god marginal inne i \code{STATE_SOF} vid det laget, och en vanlig
kombinatorisk enable räcker.

% ------------------------------------------------------------------------------------------
\section{Ett rent bord mellan ramar}
\label{c17:sec:cleanslate}

Följdspårarna nollställs bara av sin reset (\chapref{c14:ch}, \ref{c15:ch}), och inom en och samma
ram är den beständigheten precis rätt. Låt den sträcka sig \emph{över} ramgränser, och en subtil
asymmetri dyker upp: en spårares tillstånd vid SOF beror nu på vilka bitar just den noden råkade
skicka genom den tidigare, och olika noder har olika historia. En nod som sände den förra ramen
avslutar den med sin \code{tx_shift_reg}-spårare hållande CRC-fältets svans; en nod som bara tog emot
håller den historien i sin \code{rx_shift_reg}-spårare i stället, och dess sändarspårare visar
fortfarande vad nu dess \emph{egen} senaste sändning lämnade, eller ingenting alls efter reset.

Låt nu två sådana noder börja sända vid samma SOF, vilket är precis vad arbitrering finns till för.
Identiska identifierarbitar kommer ut ur båda, men noden som bär på en kvarlämnad följd når fem lika
bitar \emph{tidigare} och sätter in sin stoppbit en position före den andra noden. På den positionen
sänder den en recessiv stoppbit medan den andra noden fortfarande sänder en riktig dominant bit;
wired-AND:en läser dominant, den stoppande noden sänder recessivt inne i arbitreringsfältet, och den
drar slutsatsen att den förlorat arbitreringen. Ingen tvist ägde rum: den \emph{blivande vinnaren}
går ifrån sin egen ram med \code{error} höjd. Det här är inte hypotetiskt; det är precis så
\code{can_controller_tb}s fall 2 misslyckas om spårarna lever kvar, med nod A som går in i det
fortfarande hållande svansen av fall 1:s CRC.

Riktig CAN har inget sådant fel, eftersom stoppbitsräkningen börjar om från noll vid varje SOF, och
den här designen får samma garanti med en enda signal: båda skiftregistren \textbf{hålls i reset
medan bussen är ledig}, så varje ram börjar med rena spårare på varje nod, vad var och en än gjorde
innan.

Konkret betyder det en egen, aktivt låg reset för skiftregistren, konkurrent tilldelad: låg när
\code{reset_s2_n} är låg \textbf{eller} tillståndsmaskinen står i \code{STATE_IDLE}, hög annars.
Båda skiftregisterinstanserna från \chapref{c14:ch} och \ref{c15:ch} kopplar sin resetport, den
andra i deras portlistor, till den signalen i stället för till \code{reset_s2_n}. Ingenting annat i
instansieringarna ändras.

Reseten släpper på den flank som lämnar \code{STATE_IDLE}, en hel cykel innan SOF-laddningen klockas
in i \code{tx_shift_reg}, så registren är vakna i tid inför varje ram. Lägg märke till vad det här
\emph{inte} ändrar: inom en ram nollställs aldrig något, så det kedjade omladdningsbeteende som
\chapref{c14:ch}:s test 3 låser fast (en följd som fortsätter över en gruppgräns) är orört, och
argumentet om SOF-hoppet ovan står sig bit för bit. Det rena bordet börjar vid \code{STATE_IDLE}, och
\code{STATE_IDLE} går bara att nå mellan ramar.

% ------------------------------------------------------------------------------------------
\section{Att mata \code{crc15} på mottagarsidan}
\label{c17:sec:crc}

Spegelbilden av \chapref{c16:ch}:s grindning på sändarsidan: aktivera \code{crc15} på
\code{rx_shift_reg.real_bit_valid} (som redan utesluter både ”ingen ny bit än” och varje kastad
stoppbit), matad med \code{rx_shift_reg.real_bit}, över SOF till och med CRC-fältet. Att mata tillbaka
det mottagna CRC-fältets egna bitar återför registret till exakt noll om ramen kom fram intakt
(\chapref{c13:ch}:s generera-och-sedan-kontrollera), vilket rapporteras av \code{crc_valid}.

Det finns en \code{crc15} och en enable, så det här ersätter \chapref{c16:ch}:s regel i stället för
att stå vid sidan av den. \code{role} väljer nu källa i stället för att stänga av motorn.

Det här är det enda stället i kapitlet där formen på koden är given, och skälet är bygget, inte
designen: \code{make build-project} avgör om kontrollern har en mottagarväg genom att leta efter just
de här två uttrycken i \file{can_controller.vhd} (bilaga~\ref{app:sim} beskriver kontrollen). Stavar
du dem annorlunda hoppas \code{can_controller_tb} över i stället för att köras. Skriv dem så här:

\begin{vhdlcode}
crc_enable  <= (txsr_bit_valid and not txsr_stuff and role) or
               (rxsr_real_bit_valid and not role);
crc_data_in <= txsr_tx_bit when role = '1' else rxsr_real_bit;
\end{vhdlcode}

Låt \chapref{c16:ch}:s version stå kvar, och en mottagande nod ackumulerar aldrig något, så
\code{crc_valid} läser \code{'1'} (resetvärdet) hela vägen och kontrollen i \code{STATE_CRC_LO_WAIT}
passerar för varje ram, korrumperad eller inte.

\code{crc_clear} berörs inte av något av det här och står kvar exakt som \chapref{c16:ch} skrev den,
en puls i \code{STATE_START} oavsett roll. Den spelar större roll nu än den gjorde då: det är i det
här kapitlet noden får tre sätt att gå ur en ram i förtid, och vart och ett lämnar motorn mitt i ett
meddelande den aldrig kommer att avsluta. Ingen av avbrottsvägarna nedan behöver nollställa den på
vägen ut, eftersom nästa ram nollställer den på vägen in.

% ------------------------------------------------------------------------------------------
\section{Väntan på ett extra sampel: den avgörande tajminglärdomen}
\label{c17:sec:preempt}

På sändarsidan lägger \code{tx_shift_reg.done} redan en extra bitperiod på att låta en grupps sista
bit löpa ut innan nästa grupp laddas. En mottagare reagerar på \code{rx_shift_reg.valid} utan någon
sådan inbyggd väntan, så \code{can_controller} måste lägga till den uttryckligen vid vart och ett av
RX-sidans \code{STATE_LOAD_*}-tillstånd: \textbf{vänta på ytterligare ett \code{bt_sample}} innan den
går vidare, eftersom ledningen fortfarande visar förra gruppens sista bit under hela dess avslutande
period.

Det finns en andra hake: \code{state} blir skifttillståndet först \textbf{en cykel efter} det
\code{bt_sample} som utlöste övergången, en cykel för sent för att aktivera \code{rx_shift_reg} för
just det samplet. Därför \textbf{föregriper} det kombinatoriska styrblocket \code{rxsr_enable} (och
sätter den kommande gruppens \code{bit_count}) för exakt det samplet medan det fortfarande står i
\code{STATE_LOAD_*}-tillståndet.

\paragraph{Hur det ser ut cykel för cykel}
Här är en gruppgräns på mottagarsidan, där den 8 bitar breda höga ID-gruppen lämnar över till den 4
bitar breda gruppen med låga ID plus RTR. Cyklerna är 50 MHz-klockans, så de två
\code{bt_sample}-pulserna ligger 50 cykler isär; följden av identiska cykler mellan dem är utelämnad.

\begin{textdrawing}
klockcykel     bt_sample             state  rxsr_bit_count  rxsr_enable  rxsr_valid
n                      0         STATE_ARB_HI               8            1           0
n + 1                  1         STATE_ARB_HI               8            1           0
n + 2                  0         STATE_ARB_HI               8            1           1
n + 3                  0    STATE_LOAD_ARB_LO               4            1           0
...                    0    STATE_LOAD_ARB_LO               4            1           0
n + 51                 1    STATE_LOAD_ARB_LO               4            1           0
n + 52                 0         STATE_ARB_LO               4            1           0
\end{textdrawing}

Läs de två sampelcyklerna mot varandra. Vid \textbf{n + 1} tas den åttonde biten av den höga
ID-gruppen, och \code{rxsr_valid} dyker upp en cykel senare vid \textbf{n + 2}, vilket är när
tillståndsmaskinen bestämmer sig för att gå vidare. \code{state} läser inte \code{STATE_LOAD_ARB_LO}
förrän vid \textbf{n + 3}.

Titta nu på \textbf{n + 51}, samplet som måste fånga den \emph{första} biten av den låga ID-gruppen.
\code{state} läser fortfarande \code{STATE_LOAD_ARB_LO} i det ögonblicket. Den blir inte
\code{STATE_ARB_LO} förrän vid \textbf{n + 52}, en cykel efter att biten redan kommit och gått. Det
är hela problemet på en rad: tillståndet som namnger gruppen anländer alltid en cykel senare än
samplet som startar den.

Föregripandet är det som gör \code{rxsr_bit_count = 4} och \code{rxsr_enable = '1'} sanna vid
\textbf{n + 51} ändå, genom att driva dem från \code{STATE_LOAD_ARB_LO}-armen i stället för
\code{STATE_ARB_LO}-armen. Driv dem från \code{STATE_ARB_LO} i stället, och samma gräns ser ut så
här:

\begin{textdrawing}
klockcykel     bt_sample             state  rxsr_bit_count  rxsr_enable  rxsr_valid
n + 3                  0    STATE_LOAD_ARB_LO               0            0           0
...                    0    STATE_LOAD_ARB_LO               0            0           0
n + 51                 1    STATE_LOAD_ARB_LO               0            0           0
n + 52                 0         STATE_ARB_LO               4            1           0
\end{textdrawing}

\code{rx_shift_reg} är avstängt över det enda sampel som spelade roll. Den första biten av
identifierarens låga grupp samlas aldrig in, så gruppen blir färdig en bit kort mot ledningen, varje
senare grupp är förskjuten med en, och \code{crc15} matas med en ström som aldrig stämmer. Ingenting
rapporterar något fel vid \textbf{n + 51}; felet dyker upp vid \code{STATE_CRC_LO_WAIT}, ungefär
fyrtio bitperioder senare, som en CRC som inte går till noll.

Lägg märke till att det här inte kostar något extra på ledningen. \code{STATE_LOAD_ARB_LO} förbrukar
redan förra gruppens avslutande bitperiod, vilket är varför den går vidare på \code{bt_sample} och
samlar inget eget; föregripandet ställer bara registret i rätt konfiguration innan den perioden är
slut.

\begin{warning}
Det här mönstret upprepas vid varje gruppgräns och är designens allra viktigaste lärdom: \textbf{en
signal som ser ”klar” ut en cykel för tidigt eller för sent tar inte ut sig själv över
gruppgränserna; den ackumuleras till en tyst felinriktning} mellan sändare och mottagare som först
långt senare dyker upp som ett CRC-fel eller ett stoppbitsbrott, långt ifrån sin orsak.
\end{warning}

% ------------------------------------------------------------------------------------------
\section{Detektering av förlorad arbitrering}
\label{c17:sec:arbitration}

Exakt \secref{c10:sec:arbitration}:s regel \textemdash{} \canref{4} är den utförliga versionen
\textemdash{} kontrollerad en gång per bit vid sampelpunkten enbart under \code{STATE_ARB_HI} och
\code{STATE_ARB_LO}, vilket är precis arbitreringsfältet, ID plus RTR.
Unika ID:n garanterar att arbitreringen avgörs inom den 11 bitar långa identifieraren; kontrollen
täcker harmlöst även RTR, där varje nod här sänder dominant, en bit som aldrig kan förlora.

I \code{STATE_ARB_HI} och \code{STATE_ARB_LO}, medan \code{role = '1'}, har noden förlorat när
\code{bt_sample = '1'}, \code{txsr_tx_bit = '1'} och \code{rx_bus_s2 = '0'}: den sände recessivt och
läser dominant. På den flanken:
\begin{itemize}
  \item \code{error} går till \code{'1'}.
  \item \code{tx_done} pulsar: försöket är över, vunnet eller förlorat.
  \item \code{role} går till \code{'0'} och \code{idle_bits} till 0.
  \item Tillståndsmaskinen går tillbaka till \code{STATE_IDLE}.
\end{itemize}

Att kontrollen bara gäller medan \code{role = '1'} är det som avgränsar den till en sändande nod; de
tre termerna är själva avkänningsvillkoret.

\paragraph{Varför \code{tx_done} pulsar även här, på en ram som aldrig sändes}
Det läser fel först: \code{tx_done} betyder ”ramen blev klar”, och det blev den här inte. Läs den i
stället som \emph{”sändningsförsöket är över, och porten är din igen”} \textemdash{} vilket är den
enda fråga en anropare kan agera på. \code{error} är det som säger om försöket lyckades, och
tillsammans bär de två hela utfallet: \code{tx_done} ensam betyder sänd, \code{tx_done} med
\code{error} betyder förlorad.

Utelämna den, och anroparen har ingen flank alls att vänta på. En nod som förlorar arbitreringen
återvänder till \code{STATE_IDLE} och är fullt kapabel att sända igen, men ingenting annonserar det
någonsin, så varje anropare som väntar på \code{tx_done} innan den begär om väntar för evigt. Det är
inte hypotetiskt: det är precis så \code{register_bank} (\chapref{c18:ch}) sätter sin TX-redo-bit, och
utan den här pulsen skulle en enda förlorad arbitrering \textemdash{} normalfallet på en buss med två
noder \textemdash{} lämna hela noden oförmögen att sända tills den resettades.

Lägg märke till att det är \code{rx_bus_s2}, den synkroniserade kopian, och inte porten
\code{rx_bus}. \Secref{c11:sec:sync}:s regel gäller överallt, den här jämförelsen inräknad:
ingenting i designen läser den råa pinnen. Det vore lätt att prata sig ur den här, eftersom
sampelpunkten ligger 35 tick in i en bit som stabiliserade sig nära tick 0, så nivån är med säkerhet
stabil vid det laget. Det argumentet handlar om \emph{när} värdet är stabilt på ledningen, och
metastabilitet handlar om vad en vippa gör när den samplar en signal som ingen klocka i det här
chippet styr. De två extra vipporna kostar två klocktick av femtio och köper den enda jämförelse hela
arbitreringsupplägget vilar på.

Att nollställa \code{idle_bits} är inte städning i bokföringen. Utan den landar noden tillbaka i
\code{STATE_IDLE} med \code{idle_bits} fullt medan vinnaren är mitt i sin ram, så vinnarens nästa
dominanta bit läses som ett nytt SOF, noden går in i \code{STATE_START}, och \code{STATE_START} suddar
den \code{error} den just höjde. Samma nollställning krävs vid varje annat avbrott nedan, av samma
skäl.

% ------------------------------------------------------------------------------------------
\section{Att färdigställa processerna}
\label{c17:sec:complete}

Ingenting från \chapref{c16:ch} skrivs om. Varje ställe där \chapref{c16:ch} lämnade mottagarrollen
till det här kapitlet fylls i enligt tabellen i \secref{c17:sec:states}, och reglerna nedan är de som
inte ryms i tabellen.

\subsection*{Den kombinatoriska processen}
Den får två tillägg, och \code{crc_valid} behöver in i dess känslighetslista, eftersom kvitteringen
läser den:
\begin{itemize}
  \item \textbf{Föregripandet.} När \code{role = '0'} driver den \code{rxsr_enable = '1'} och
    \code{rxsr_bit_count} enligt tabellens kolumn \emph{Föregriper}: i \code{STATE_SOF} med bredden
    1, och i varje stoppat fält i \textbf{både} \code{STATE_LOAD_*}-tillståndet och dess
    skifttillstånd. I alla andra tillstånd gäller förvalen från \chapref{c16:ch}. Att paret delar
    samma värde är det som gör föregripandet automatiskt: det \code{bt_sample} som faller medan
    \code{state} fortfarande läser \code{STATE_LOAD_*} samlas in med den kommande gruppens bredd,
    precis som \secref{c17:sec:preempt} härledde.
  \item \textbf{Kvitteringen.} I \code{STATE_ACK_SLOT}, när \code{role = '0'} och
    \code{crc_valid = '1'}: \code{bus_en = '1'} och \code{tx_bus = '0'}. Det är den enda bit en
    mottagare någonsin driver.
\end{itemize}

\subsection*{Den sekventiella processen}
Mottagarhalvan av varje tillstånd följer en av två former, och tabellen säger vilken:
\begin{itemize}
  \item Ett mottagande \code{STATE_LOAD_*}-tillstånd går vidare till sitt skifttillstånd på nästa
    \code{bt_sample}, och gör inget annat.
  \item Ett mottagande skifttillstånd går vidare på \code{rxsr_valid} och skördar sin grupp enligt
    tabellens sista kolumn. \code{STATE_CRC_HI} och \code{STATE_CRC_LO} skördar ingenting: deras
    bitar når \code{crc15} en i taget genom \code{real_bit}, och de hopsatta grupperna betyder
    ingenting.
\end{itemize}

Tre regler till, som tabellen inte uttrycker:
\begin{itemize}
  \item \textbf{Arbitreringskontrollen och \code{done}.} I \code{STATE_ARB_HI} och
    \code{STATE_ARB_LO} har sändarhalvan nu två villkor: förlorad arbitrering enligt
    \secref{c17:sec:arbitration}, och \code{txsr_done} som i \chapref{c16:ch}. De två kan aldrig slå
    till samma cykel, eftersom \code{bt_sample} ligger mitt i perioden och \code{done} följer ett
    \code{bit_done}, så ordningen handlar om läsbarhet snarare än prioritet; en förlust på en grupps
    sista bit vinner ändå helt enkelt genom att komma först, vid den bitens sampelpunkt.
  \item \textbf{Två läsningar före flanken}, båda samma regel som \chapref{c14:ch}:s
    \code{MAX_RUN - 1}. \code{STATE_CTRL} avgör DLC-grenen direkt utifrån
    \code{rxsr_data(3 downto 0)}, eftersom \code{rx_dlc_acc} tilldelas på just den flanken och
    fortfarande läser den förra ramens värde. Och \code{STATE_DATA} jämför mot \code{rx_dlc_acc},
    som vid det laget \emph{är} stabil, låst ett helt fält tidigare.
  \item \textbf{Ett stoppbitsbrott avbryter från vilket tillstånd som helst.} När \code{role = '0'}
    och \code{rxsr_stuff_error = '1'}: \code{error} till \code{'1'}, \code{idle_bits} till 0, och
    tillbaka till \code{STATE_IDLE}. Regeln måste vinna över vad tillståndets egen regel bestämde på
    samma flank, vilket är hela poängen: ett brott är fatalt var det än slår till, och ett sampel kan
    slå till i vilket som helst av tillstånden för stoppade fält, LOAD-tillstånden inräknade (deras
    föregripna sampel behandlas som vilket annat som helst). I en klockad process är det den sista
    tilldelningen till en signal som gäller, så att skriva regeln en gång, efter alla tillstånd, är
    bättre än att upprepa den i ett dussin.
\end{itemize}

\subsection*{Att kvittera, kontrollera och avbryta}
\begin{itemize}
  \item \textbf{ACK-luckan:} en mottagare vars \code{crc_valid} är satt driver bussen dominant under
    den enda biten för att kvittera ramen.
  \item \textbf{CRC-kontrollen:} i \code{STATE_CRC_LO_WAIT} läser en mottagare \code{crc_valid}. Är
    den satt väntar den på \code{bt_bit_done} och går sedan vidare till \code{STATE_CRC_DELIM}; är
    den det inte avbryter den.
  \item \textbf{Stoppbitsbrott:} \code{rx_shift_reg.stuff_error} avbryter mottagningen från vilket
    tillstånd som helst. En sändare kan inte bryta mot stoppbitsregeln, så det här gäller bara RX.
  \item \textbf{Arbitreringsförlust} är det enda avbrott som sker under \emph{sändning}, och därmed
    det enda som också pulsar \code{tx_done}. Avbrotten på mottagarsidan avslutar en mottagning, och
    en mottagning har ingen anropare som väntar på en färdigpuls: \code{rx_valid} uteblir helt
    enkelt, vilket är precis rätt, eftersom ingen ram anlände.
  \item \textbf{\code{error} ligger kvar:} nollställd i \code{STATE_START}, hållen i övrigt, så att
    en anropare ser den ända tills nästa ram börjar. Det är precis därför \code{STATE_IDLE} kräver en
    \textbf{ledig buss} innan noden lämnar det tillståndet över huvud taget, och varför varje avbrott
    här måste nollställa \code{idle_bits} på vägen ut. Varje avbrott landar tillbaka i
    \code{STATE_IDLE} mitt i en ram, medan bussen fortfarande är upptagen. Utan kravet på ledig buss
    skulle noden gå in i \code{STATE_START} igen inom några få bitperioder och sudda den
    \code{error} den höjde, så att en anropare som pollar \code{error} aldrig skulle se den. Att
    först kräva elva recessiva bitperioder är exakt den garanti som ”ända tills nästa ram börjar”
    utfärdar.
\end{itemize}

% ------------------------------------------------------------------------------------------
\section{Att verifiera den: \code{can_controller_tb}}
\label{c17:sec:verify}

Nu när mottagarvägen är komplett körs \file{controller/can_controller_tb.vhd} äntligen hela vägen.
Den kopplar \textbf{två} \code{can_controller}-noder till en enda buss med öppen dränering, kombinerar
varje nods par \code{(tx_bus, bus_en)} på det sätt \secref{c10:sec:can} beskrev en wired-AND, och kör
två fall.

\paragraph{Fall 1, mottagning}
Nod A sänder ID \code{0x123}, DLC 5, data \code{11 22 F8 44 55} medan nod B tar emot. Den
kontrollerar \code{tx_done} och \code{rx_valid}, och att B återskapar ID:t, DLC:n och \textbf{varje}
databyte på sin fasta position. Byten \code{0xF8} är vald för att tvinga fram en stoppbit inne i
datafältet, så det här övar sändvägen, mottagarvägen, avstoppningen och flerbytesplaceringen i en och
samma ram. En ensam nod kan inte testa något av det: \code{rx_shift_reg} är aktivt bara medan
\code{role = '0'}, så en sändare tar aldrig emot sin egen ram.

\paragraph{Fall 2, arbitrering}
Båda noderna begär under samma cykel, A med ID \code{0x000} och B med \code{0x001}, åtskilda bara i
den sista identifierarbiten. B måste upptäcka förlusten på exakt den biten och sluta driva bussen; A
måste bli klar som vanligt, ovetande om att en tvist ägt rum. Den sista assertionen är den att lägga
märke till: den läser B:s \code{error} \emph{efter} att A:s ram blivit klar, så den faller om inte
både den kvarliggande \code{error} och disciplinen kring \code{idle_bits} ovan håller.
